% =====================================================================
% Guía del estudiante -- Geometría 10° -- Trimestre I
% María Mercedes Cantillo · Colegio San Alberto Magno
% ---------------------------------------------------------------------
% Hilo conductor: Ubicarse en la Tierra (de las coordenadas geográficas
% al plano cartesiano)
% Plan (etapas A y B): guia-periodo-I-geometria-analitica.plan.md
% Temas (propuesta de curriculo/plan-anual.md, sesiones 001-009):
%   planocartesiano, distancia, lugaresgeometricos, recta
% Cada \tema tiene \label{tema:...} para que las clases lo citen con
% \refguia{tema:...}. No borrar el .aux de esta guía.
% Ejercicios (reutilizados del banco salvo 4 nuevos): geometria-analitica-10,
% plano-cartesiano-5, coordenadas-11, lugares-geometricos-11, funcion-lineal-9
% e icfes-cuadernillo-2026 (textuales).
%
% Compilar:  latexmk -pdf guia-periodo-I-geometria-analitica.tex
% =====================================================================
% BORRADOR: el plan del trimestre está pendiente de aprobación de la docente (generación rápida 2026-09-14).
\documentclass[11pt]{article}
\makeatletter\def\input@path{{../../../../plantillas/estilo/}}\makeatother
\usepackage{mmcantillo}

% ======================== DATOS DE LA GUÍA ============================
\tipodocumento{Guía del estudiante}
\asignatura{Geometría}
\titulo{El plano cartesiano y la recta}
\grado{10°}
\periodo{I}
\hiloconductor{Ubicarse en la Tierra: de las coordenadas geográficas al plano cartesiano}
% ======================================================================
% DBA: matematicas grado 10 · DBA 5 — Explora y describe las propiedades de
%      los lugares geométricos y de sus transformaciones a partir de
%      diferentes representaciones.

\begin{document}

\encabezado

% ---------------------------------------------------------------------
% 1. Frase introductoria
% ---------------------------------------------------------------------
% verificado: https://fr.wikisource.org/wiki/La_G%C3%A9om%C3%A9trie_(%C3%A9d._1637) — primera frase del libro primero: «Tous les Problesmes de Geometrie se peuuent facilement reduire a tels termes, qu'il n'est besoin par aprés que de connoitre la longueur de quelques lignes droites, pour les construire» (Leiden, Jan Maire, 1637)
\frase{Todos los problemas de geometría pueden reducirse fácilmente a tales
términos que después no es necesario conocer más que la longitud de algunas
líneas rectas para construirlos.}
{René Descartes, \emph{La geometría} (1637), libro primero (trad. propia)}

% ---------------------------------------------------------------------
% 2. Introducción
% ---------------------------------------------------------------------
\seccion{Introducción}

En la primera clase vimos cómo se ubica cualquier lugar de la Tierra con dos
números: la latitud y la longitud. Un mapa de Cali hace lo mismo con una
cuadrícula: basta decir «tantas cuadras al oriente y tantas al norte». Esa idea
—dos rectas perpendiculares y dos números— es el \textbf{plano cartesiano}, y
con él la geometría se vuelve álgebra: una distancia se calcula con una
fórmula, y una figura se describe con una ecuación.

Nuestro hilo conductor es \emph{ubicarse en la Tierra}: de Hiparco, que usó la
latitud y la longitud para fijar lugares, a Descartes y Fermat, que llevaron las
coordenadas a la geometría, y a Katherine Johnson, que calculó con ecuaciones
dónde debía caer una nave espacial. Cada tema responde una pregunta de
ubicación: ¿dónde está?, ¿a qué distancia?, ¿qué lugares cumplen una
condición?, ¿cómo describo un camino recto?

\textbf{Al terminar este trimestre podrás:}
\begin{itemize}
  \item ubicar puntos y figuras en el plano cartesiano y relacionarlo con las
        coordenadas geográficas;
  \item calcular la distancia entre dos puntos y el punto medio de un segmento,
        y usarlos para clasificar figuras;
  \item describir la mediatriz y la circunferencia como lugares geométricos y
        escribir sus ecuaciones;
  \item hallar la pendiente y la ecuación de una recta, graficarla y reconocer
        rectas paralelas y perpendiculares.
\end{itemize}

Este trimestre trabaja un Derecho Básico de Aprendizaje del Ministerio de
Educación Nacional~\cite{mendba}:

\begin{dba}[Grado 10° · DBA 5]
Explora y describe las propiedades de los lugares geométricos y de sus
transformaciones a partir de diferentes representaciones.
\end{dba}

% ---------------------------------------------------------------------
% 3. Marco teórico
% ---------------------------------------------------------------------
\seccion{Marco teórico}

\subsubsection*{Dos números para un lugar}

Hiparco de Nicea (c.~190--120 a.\,C.)~\cite{mactutorhiparco}, el gran
astrónomo griego, comenzó a usar de manera sistemática la latitud y la
longitud —que ya empleaba en su catálogo de estrellas— para fijar lugares
sobre la superficie de la Tierra~\cite{ebsco}. La \textbf{latitud} dice qué
tan al norte o al sur del ecuador está un lugar; la \textbf{longitud}, qué tan
al oriente o al occidente de un meridiano de referencia. Hoy, por ejemplo, la
Alcaldía ubica a Cali en $3^\circ 27' 00''$ de latitud norte y
$76^\circ 32' 00''$ de longitud oeste~\cite{alcaldia}.

\subsubsection*{Del globo al plano}

La Tierra es casi una esfera, pero en el mapa de una ciudad o de una región
pequeña podemos tratarla como si fuera plana y dibujarle una cuadrícula. Cada
grado de latitud equivale a unos $111$ km sobre la superficie
terrestre~\cite{noaa}; en distancias cortas, el error de «aplanar» la Tierra
es pequeño. Esa cuadrícula, con dos ejes perpendiculares, es el plano
cartesiano.

\subsubsection*{Descartes y Fermat: la geometría se vuelve álgebra}

En 1637, René Descartes publicó en Leiden su \emph{Discurso del método} con
tres apéndices; el más importante fue \emph{La Géométrie}, donde aplicó el
álgebra a la geometría~\cite{mactutordescartes,descartes}. Por la misma época,
y de forma independiente, Pierre de Fermat había llegado a un sistema casi
idéntico: su \emph{Introducción a los lugares planos y sólidos} ya estaba
escrita en la primavera de 1636 y circuló en manuscrito. Allí la ecuación más
sencilla, $Dx = Ey$, representa una recta~\cite{mahoney}. La obra se publicó
después de su muerte, en 1679~\cite{fermat}. De Fermat nos queda la palabra
\emph{lugar}: un lugar geométrico es el conjunto de todos los puntos que
cumplen una condición.

\subsubsection*{Katherine Johnson: ecuaciones para aterrizar}

Katherine Johnson (1918--2020) fue matemática de la NASA. En 1960 escribió con
el ingeniero Ted Skopinski un informe con las ecuaciones de un vuelo orbital
en el que se especifica el lugar de aterrizaje de la nave~\cite{skopinski}. Hizo
el análisis de la trayectoria del vuelo de Alan Shepard en 1961 y, en 1962,
comprobó a mano las ecuaciones orbitales programadas en los computadores para
el vuelo de John Glenn. En 2015 recibió la Medalla Presidencial de la
Libertad~\cite{nasajohnson}. Su trabajo fue, literalmente, ubicar puntos sobre la
Tierra con ecuaciones.

% ---------------------------------------------------------------------
% 4. Temas
% ---------------------------------------------------------------------
\tema{Plano cartesiano}\label{tema:planocartesiano}

\begin{hilo}
Hiparco necesitó dos números para ubicar una ciudad: latitud y longitud. En un
mapa pequeño, los meridianos y los paralelos se ven casi como rectas
perpendiculares: una cuadrícula. Ese es nuestro punto de partida.
\end{hilo}

\begin{definicion}[Plano cartesiano]
Dos rectas numéricas perpendiculares, el \textbf{eje $x$} (horizontal) y el
\textbf{eje $y$} (vertical), que se cortan en el \textbf{origen} $O(0, 0)$.
Cada punto $P$ del plano se nombra con una \textbf{pareja ordenada}
$(x, y)$: primero la abscisa $x$ (cuánto a la derecha o a la izquierda del
origen) y después la ordenada $y$ (cuánto arriba o abajo). Los ejes dividen el
plano en cuatro \textbf{cuadrantes}, numerados I, II, III y IV en sentido
contrario a las manecillas del reloj; los puntos sobre los ejes no están en
ningún cuadrante.
\end{definicion}

% Ejercicio: geometria-analitica-10-002
\begin{ejemplo}[Cuadrantes y ejes]
¿Dónde está cada punto? $(-3, 5)$ está en el cuadrante II; $(4, -2)$, en el
IV; $(-1, -1)$, en el III. $(0, -6)$ está sobre el eje $y$ y $(7, 0)$ sobre el
eje $x$: tienen una coordenada igual a $0$. Los signos dicen el cuadrante:
$(+,+)$ I, $(-,+)$ II, $(-,-)$ III, $(+,-)$ IV.

\begin{center}
\begin{tikzpicture}
\begin{axis}[mmc, xmin=-4.5, xmax=8, ymin=-7, ymax=6,
    xtick={-3,...,7}, ytick={-6,...,5}, width=0.55\linewidth, height=0.5\linewidth]
  \addplot[only marks, mark=*, mark size=1.8pt] coordinates {(-3,5) (4,-2) (0,-6) (-1,-1) (7,0)};
  \node[above, fill=white, inner sep=1pt] at (axis cs:-3,5.2) {$(-3, 5)$};
  \node[right] at (axis cs:4,-2) {$(4, -2)$};
  \node[right] at (axis cs:0,-6) {$(0, -6)$};
  \node[below, fill=white, inner sep=1pt] at (axis cs:-1.8,-1.3) {$(-1, -1)$};
  \node[above] at (axis cs:7,0) {$(7, 0)$};
  \node[font=\small\bfseries] at (axis cs:3,4) {I};
  \node[font=\small\bfseries] at (axis cs:-3,2.5) {II};
  \node[font=\small\bfseries] at (axis cs:-3,-4) {III};
  \node[font=\small\bfseries] at (axis cs:5,-4.5) {IV};
\end{axis}
\end{tikzpicture}
\end{center}
\end{ejemplo}

Con puntos se forman figuras: un polígono queda determinado por las
coordenadas de sus vértices. Y un mismo lugar puede tener coordenadas distintas
si cambiamos el origen: las coordenadas dependen del \emph{sistema de
referencia}, pero las distancias no.

\begin{aplicacion}[Cali en el plano]
Si tomamos la longitud como eje $x$ (positiva hacia el este) y la latitud como
eje $y$ (positiva hacia el norte), Cali ($3^\circ 27'$ N, $76^\circ 32'$ O) es
un punto con $x$ negativa e $y$ positiva: está en el «cuadrante II» del mundo,
al oeste del meridiano de Greenwich y al norte del ecuador. Por eso muchos
mapas digitales escriben su ubicación con signos, como $(3{,}45;\ -76{,}53)$:
latitud primero y longitud después, con negativo para el oeste.
\end{aplicacion}

\begin{preguntas}
% Ejercicio: plano-cartesiano-5-013
\Needspace{6\baselineskip}
\pregunta En un plano cartesiano están marcados los puntos $D(-1, +1)$,
$H(+2, +1)$, $K(+5, +1)$, $A(-5, +2)$, $C(-3, +2)$, $E(0, +2)$, $B(-3, +4)$,
$F(0, +4)$, $G(+2, +4)$, $J(+5, +4)$, $N(-3, -2)$, $L(+2, -2)$, $M(+4, -2)$,
$P(-3, -4)$. ¿Cuál es mayor: la distancia del punto $B$ al punto $N$ o la
distancia del punto $A$ al punto $E$?
\espacio{3cm}
% Ejercicio: coordenadas-11-003
\pregunta En un mapa con unidades en kilómetros y origen en el puerto, una
estación de radio está en $E(0, 0)$ y un barco en $B(30, 40)$. Un segundo mapa
pone el origen en un faro que está $20$~km al este y $10$~km al norte del
puerto.
\begin{opciones}
  \item Escribe las coordenadas de $E$ y $B$ en el segundo mapa.
  \item Calcula la distancia $EB$ en los dos mapas. ¿Qué cambia y qué no cambia
        al mover el origen?
\end{opciones}
\espacio{3.5cm}
% Ejercicio: geometria-analitica-10-006
\pregunta Según la Alcaldía, Cali está en $3^\circ 27'$ de latitud norte y
$76^\circ 32'$ de longitud oeste. Toma un plano en el que $x$ es la longitud
(positiva al este) e $y$ la latitud (positiva al norte), en grados. Escribe la
ubicación de Cali como un par ordenado con dos decimales y di en qué cuadrante
queda. (Recuerda: $1^\circ = 60'$.)
\espacio{2.5cm}
\end{preguntas}

\begin{resumen}
\begin{itemize}
  \item Un punto del plano es una pareja \emph{ordenada} $(x, y)$: el orden
        importa.
  \item Los signos de las coordenadas indican el cuadrante; si una coordenada
        es $0$, el punto está sobre un eje.
  \item Al cambiar el origen cambian las coordenadas, no las distancias.
        Latitud y longitud son coordenadas sobre la esfera; en un mapa pequeño
        funcionan como un plano cartesiano.
\end{itemize}
\end{resumen}

% ---------------------------------------------------------------------
\tema{Distancia entre dos puntos}\label{tema:distancia}

\begin{hilo}
Ya sabemos \emph{dónde} está cada lugar. La siguiente pregunta de quien mira un
mapa es \emph{¿a qué distancia?}, y si dos personas quieren encontrarse a mitad
de camino, \emph{¿dónde?}
\end{hilo}

Entre $P_1(x_1, y_1)$ y $P_2(x_2, y_2)$ se forma un triángulo rectángulo con
catetos $|x_2 - x_1|$ y $|y_2 - y_1|$. Por el teorema de Pitágoras, la
hipotenusa es la distancia.

\begin{teorema}[Distancia y punto medio]
\[
  d(P_1, P_2) = \sqrt{(x_2 - x_1)^2 + (y_2 - y_1)^2},
  \qquad
  M = \left( \frac{x_1 + x_2}{2},\ \frac{y_1 + y_2}{2} \right).
\]
El punto medio $M$ del segmento $\overline{P_1P_2}$ es el promedio de las
coordenadas.
\end{teorema}

% Ejercicio: coordenadas-11-001
\begin{ejemplo}[Distancia y punto medio]
Entre $A(-3, 2)$ y $B(5, 8)$ los catetos miden $5 - (-3) = 8$ y $8 - 2 = 6$:
\[ d = \sqrt{8^2 + 6^2} = \sqrt{100} = 10, \qquad
   M = \left(\frac{-3 + 5}{2}, \frac{2 + 8}{2}\right) = (1, 5). \]

\begin{center}
\begin{tikzpicture}
\begin{axis}[mmc, xmin=-4, xmax=6.5, ymin=0, ymax=9.5, xtick={-3,...,5}, ytick={2,4,...,8},
    width=0.5\linewidth, height=0.42\linewidth]
  \addplot[black, very thick] coordinates {(-3,2) (5,8)};
  \addplot[black, thick, dashed, sharp plot] coordinates {(-3,2) (5,2) (5,8)};
  \addplot[only marks, mark=*, mark size=1.8pt] coordinates {(-3,2) (5,8) (1,5)};
  \node[below left] at (axis cs:-3,2) {$A$};
  \node[above left] at (axis cs:5,8) {$B$};
  \node[above left] at (axis cs:1,5) {$M$};
  \node[below] at (axis cs:1,2) {$8$};
  \node[right] at (axis cs:5,5) {$6$};
\end{axis}
\end{tikzpicture}
\end{center}
\end{ejemplo}

Con la distancia podemos \textbf{clasificar triángulos} dados sus vértices:
por sus lados (equilátero, isósceles, escaleno) y por sus ángulos, usando el
recíproco del teorema de Pitágoras: si $a^2 + b^2 = c^2$ para los tres lados,
el triángulo es rectángulo.

% Ejercicio: geometria-analitica-10-012
\begin{aplicacion}[Distancias en el globo]
Cada grado de latitud equivale a unos $111$ km~\cite{noaa}. Dos lugares sobre
el mismo meridiano, uno en $2^\circ 15'$ N y otro en $3^\circ 45'$ N, están
separados por $1^\circ 30' = \num{1,5}^\circ$ de latitud, es decir, por unos
$\num{1,5} \cdot 111 \approx \num{166,5}$ km en dirección norte--sur. Para
distancias en otras direcciones se necesita la geometría de la esfera, pero en
el mapa de una ciudad basta la fórmula de la distancia.
\end{aplicacion}

\begin{preguntas}
% Ejercicio: coordenadas-11-002
\pregunta Halla la distancia y el punto medio de cada par de puntos:
\begin{opciones*}(3)
  \item $(1, 1)$ y $(4, 5)$
  \item $(-2, 3)$ y $(4, -5)$
  \item $(0, 0)$ y $(3, 3)$
\end{opciones*}
\espacio{3cm}
% Ejercicio: coordenadas-11-004
\Needspace{14\baselineskip}
\pregunta Demuestra con coordenadas que el triángulo de vértices $A(1, 1)$,
$B(5, 4)$ y $C(2, 8)$ es isósceles y rectángulo.
\espacio{3.5cm}
% Ejercicio: coordenadas-11-005
\pregunta Un dron de rescate sale del punto $(2, -1)$ y debe llegar a
$(14, 4)$ (coordenadas en km).
\begin{opciones}
  \item ¿Qué distancia recorre en línea recta?
  \item Debe reportarse a mitad de camino: ¿en qué punto?
  \item Si vuela a $52$~km/h, ¿cuántos minutos tarda?
\end{opciones}
\espacio{3.5cm}
\end{preguntas}

\begin{resumen}
\begin{itemize}
  \item La distancia entre dos puntos es la hipotenusa del triángulo cuyos
        catetos son las diferencias de las coordenadas.
  \item El punto medio es el promedio de las coordenadas de los extremos.
  \item Con las distancias se clasifican triángulos; con el recíproco de
        Pitágoras se reconoce el ángulo recto.
\end{itemize}
\end{resumen}

% ---------------------------------------------------------------------
\tema{Lugares geométricos}\label{tema:lugaresgeometricos}

\begin{hilo}
Fermat hablaba de \emph{lugares}: no un punto, sino todos los puntos que
cumplen una condición. ¿Dónde tiene señal una antena? ¿Dónde poner algo que
quede igual de lejos de dos estaciones? Las respuestas son lugares
geométricos.
\end{hilo}

\begin{definicion}[Lugar geométrico]
Un \textbf{lugar geométrico} es el conjunto de todos los puntos del plano que
cumplen una condición. Su \textbf{ecuación} es la relación entre $x$ e $y$ que
satisfacen esos puntos, y solo ellos.
\end{definicion}

\begin{itemize}
  \item La \textbf{mediatriz} de un segmento $\overline{AB}$ es el lugar de los
        puntos $P$ que están a igual distancia de $A$ y de $B$: $PA = PB$. Es la
        recta perpendicular a $\overline{AB}$ por su punto medio.
  \item La \textbf{circunferencia} de centro $C(h, k)$ y radio $r$ es el lugar
        de los puntos que están a distancia $r$ de $C$:
        $(x - h)^2 + (y - k)^2 = r^2$.
\end{itemize}

% Ejercicio: lugares-geometricos-11-004
\begin{ejemplo}[La mediatriz]
Dos estaciones están en $A(0, 0)$ y $B(8, 4)$. Los puntos $P(x, y)$ a igual
distancia de las dos cumplen $PA^2 = PB^2$:
\[ x^2 + y^2 = (x - 8)^2 + (y - 4)^2. \]
Al desarrollar, $x^2$ e $y^2$ se cancelan: $0 = -16x + 64 - 8y + 16$, es decir,
$16x + 8y = 80$, o $y = -2x + 10$. Pasa por el punto medio $(4, 2)$, y el
punto $(5, 0)$ la cumple: está a $5$ de $A$ y a $\sqrt{9 + 16} = 5$ de $B$.
\end{ejemplo}

% Ejercicio: lugares-geometricos-11-001
\begin{ejemplo}[La circunferencia]
La circunferencia de centro $(3, -2)$ y radio $5$ es
$(x - 3)^2 + (y + 2)^2 = 25$. El punto $(6, 2)$ está sobre ella, porque
$(6 - 3)^2 + (2 + 2)^2 = 9 + 16 = 25$.
\end{ejemplo}

\begin{center}
\begin{minipage}{\linewidth}\centering
\begin{tikzpicture}
\begin{axis}[mmc, axis equal, xmin=-3, xmax=10, ymin=-8, ymax=11,
    xtick={-2,2,4,6,8}, ytick={-6,-4,-2,2,4,6,8,10}, width=0.55\linewidth, height=0.6\linewidth]
  \addplot+[domain=0:360, samples=120, variable=\t] ({3+5*cos(t)}, {-2+5*sin(t)});
  \addplot+[domain=-0.5:5.5] {10-2*x};
  \addplot[black, thick] coordinates {(0,0) (8,4)};
  \addplot[only marks, mark=*, mark size=1.6pt] coordinates {(3,-2) (6,2) (0,0) (8,4) (4,2)};
  \node[below right] at (axis cs:3,-2) {$(3, -2)$};
  \node[below right, fill=white, inner sep=1pt] at (axis cs:6.1,1.9) {$(6, 2)$};
  \node[left] at (axis cs:0,0) {$A$};
  \node[above] at (axis cs:8,4) {$B$};
\end{axis}
\end{tikzpicture}

{\small Circunferencia de centro $(3, -2)$ y radio $5$ (línea continua) y
mediatriz $y = -2x + 10$ del segmento $\overline{AB}$ (línea discontinua).}
\end{minipage}
\end{center}

\begin{aplicacion}[La cobertura de una antena]
Si en terreno plano y sin obstáculos una antena llega hasta $3$ km a la
redonda, la zona con señal es el círculo de radio $3$ con centro en la antena,
y su borde es una circunferencia. En la realidad los edificios y las montañas
deforman esa zona, pero el modelo ideal ya responde la pregunta básica: una casa
tiene señal si su distancia a la antena es menor o igual que $3$ km.
\end{aplicacion}

\begin{preguntas}
% Ejercicio: lugares-geometricos-11-003
\Needspace{9\baselineskip}
\pregunta Una estación en el origen emite un pulso de radio que tarda
$200$~µs en llegar a un barco. Las ondas de radio viajan a unos
$\num{0,3}$~km/µs.
\begin{opciones}
  \item ¿A qué distancia está el barco?
  \item ¿Dónde puede estar? Escribe la ecuación.
  \item ¿Por qué una sola estación no basta para saber la posición?
\end{opciones}
\espacio{3.5cm}
% Ejercicio: lugares-geometricos-11-005
\pregunta Pedro dice: «Los puntos que están a igual distancia de $A(1, 2)$ y
$B(5, 2)$ forman la circunferencia de centro $(3, 2)$ que pasa por $A$ y
$B$». Da un punto de esa circunferencia que no esté a igual distancia de $A$ y
$B$, y escribe el lugar geométrico correcto.
\espacio{3.5cm}
\end{preguntas}

\begin{resumen}
\begin{itemize}
  \item Un lugar geométrico es el conjunto de \emph{todos} los puntos que
        cumplen una condición; su ecuación la traduce al álgebra.
  \item Mediatriz: $PA = PB$; es una recta perpendicular al segmento por su
        punto medio.
  \item Circunferencia de centro $(h, k)$ y radio $r$:
        $(x - h)^2 + (y - k)^2 = r^2$.
\end{itemize}
\end{resumen}

% ---------------------------------------------------------------------
\tema{La recta}\label{tema:recta}

\begin{hilo}
Katherine Johnson calculaba trayectorias de naves que daban vueltas a la
Tierra. Antes de las curvas, el camino más sencillo entre dos lugares de un
mapa es la línea recta, el primer lugar geométrico que Fermat escribió con una
ecuación.
\end{hilo}

\begin{definicion}[Pendiente]
La \textbf{pendiente} de la recta que pasa por $P_1(x_1, y_1)$ y
$P_2(x_2, y_2)$, con $x_1 \neq x_2$, es
\[ m = \frac{y_2 - y_1}{x_2 - x_1}, \]
es decir, cuánto cambia $y$ por cada unidad que avanza $x$. Si $m > 0$ la
recta es creciente; si $m < 0$, decreciente; si $m = 0$, horizontal. Si
$x_1 = x_2$, la recta es vertical y su pendiente no está definida.
\end{definicion}

% Ejercicio: funcion-lineal-9-046
\begin{ejemplo}[Calcular una pendiente]
Por $(2, 5)$ y $(6, 10)$: $m = \frac{10 - 5}{6 - 2} = \frac{5}{4}$. Por cada
$4$ unidades que avanza $x$, $y$ sube $5$.
\end{ejemplo}

\begin{teorema}[Ecuaciones de la recta]
\begin{itemize}
  \item Punto-pendiente: $y - y_1 = m(x - x_1)$.
  \item Pendiente-intercepto: $y = mx + b$, donde $(0, b)$ es el corte con el
        eje $y$.
  \item General: $Ax + By + C = 0$, con $A$ y $B$ no ambos cero. Si $B \neq 0$,
        $m = -\frac{A}{B}$.
\end{itemize}
Dos rectas no verticales son \textbf{paralelas} si $m_1 = m_2$, y
\textbf{perpendiculares} si $m_1 \cdot m_2 = -1$. Una recta horizontal y una
vertical también son perpendiculares.
\end{teorema}

% Ejercicio: funcion-lineal-9-060
\begin{ejemplo}[La recta que pasa por dos puntos]
Por $(5, 3)$ y $(2, 8)$: $m = \frac{8 - 3}{2 - 5} = -\frac{5}{3}$. Con el punto
$(5, 3)$: $y - 3 = -\frac{5}{3}(x - 5)$, y despejando,
$y = -\frac{5}{3}x + \frac{34}{3}$.
\end{ejemplo}

% Ejercicio: funcion-lineal-9-011
\begin{ejemplo}[Graficar desde la ecuación]
Para trazar $y = -2x + 4$ bastan dos puntos: con $x = 0$, $y = 4$; con $y = 0$,
$x = 2$. La recta corta los ejes en $(0, 4)$ y $(2, 0)$ y es decreciente
($m = -2$).

\begin{center}
\begin{tikzpicture}
\begin{axis}[mmc, xmin=-1.5, xmax=4, ymin=-2.5, ymax=6, xtick={-1,...,3}, ytick={-2,...,5},
    width=0.45\linewidth, height=0.42\linewidth]
  \addplot+[domain=-1:3.2] {-2*x+4};
  \addplot[only marks, mark=*, mark size=1.8pt] coordinates {(0,4) (2,0)};
  \node[right] at (axis cs:0,4) {$(0, 4)$};
  \node[above right] at (axis cs:2,0) {$(2, 0)$};
\end{axis}
\end{tikzpicture}
\end{center}
\end{ejemplo}

\begin{aplicacion}[Una ruta recta en el mapa]
Si la ruta de un bus en el mapa es un tramo recto, su ecuación permite saber,
sin dibujar, si pasa por un lugar: basta reemplazar las coordenadas del lugar
en la ecuación. Si la igualdad se cumple, el lugar está sobre la ruta; si no,
queda por fuera. Y si dos rutas tienen la misma pendiente y distinto corte con
el eje $y$, son paralelas: nunca se cruzan.
\end{aplicacion}

\begin{preguntas}
% Ejercicio: funcion-lineal-9-061
\pregunta Encuentra la ecuación de la recta que pasa por $(4, -4)$ y
$(-8, 8)$.
\espacio{2.5cm}
% Ejercicio: funcion-lineal-9-069
\Needspace{8\baselineskip}
\pregunta Encuentra la ecuación de una recta paralela y la de una recta
perpendicular a la recta $y = -\frac{5}{3}x + \frac{34}{3}$ (la que pasa por
$(5, 3)$ y $(2, 8)$).
\espacio{3cm}
% Ejercicio: funcion-lineal-9-094
\pregunta Grafica las rectas $y = 5$ y $x = 4$ y determina si son paralelas,
perpendiculares o ninguna de las dos.
\espacio{3cm}
% Ejercicio: funcion-lineal-9-095
\pregunta Grafica las rectas $2x + 3y = -2$ y $3x - y = 4$ y determina si son
paralelas, perpendiculares o ninguna de las dos.
\espacio{3cm}
% Ejercicio: geometria-analitica-10-030
\pregunta Sofía calculó la pendiente de la recta que pasa por $(1, 3)$ y
$(4, 9)$ así: $m = \frac{4 - 1}{9 - 3} = \frac{1}{2}$. Encuentra el error y
corrígelo.
\espacio{2.5cm}
\end{preguntas}

\begin{resumen}
\begin{itemize}
  \item Pendiente: $m = \frac{y_2 - y_1}{x_2 - x_1}$; su signo dice si la
        recta crece o decrece; la vertical no tiene pendiente.
  \item Formas de la ecuación: $y - y_1 = m(x - x_1)$, $y = mx + b$ y
        $Ax + By + C = 0$.
  \item Paralelas: $m_1 = m_2$. Perpendiculares: $m_1 \cdot m_2 = -1$.
\end{itemize}
\end{resumen}

% ---------------------------------------------------------------------
% 5. Prepárate para Saber 11
% ---------------------------------------------------------------------
\seccion{Prepárate para Saber 11}

\begin{instrucciones}
Preguntas de selección múltiple con única respuesta, tomadas textualmente del
cuadernillo de preguntas de Matemáticas Saber 11.º del Icfes~\cite{icfes}.
Resuélvelas en tu cuaderno y compara tus respuestas con las de tus compañeros.
\end{instrucciones}

\begin{preguntas}
% Ejercicio: icfes-cuadernillo-2026-008 (textual; © Icfes 2026, uso académico)
\pregunta Un grupo de montañistas sabe que cada vez que aumenta la altitud en
100 m, la temperatura disminuye en 1~°C. Si el grupo se encuentra a una
altitud de 1.000 m, donde la temperatura es de 20~°C, ¿cuál de las siguientes
expresiones les permite determinar la temperatura que habrá cuando se
encuentren a 4.000 m de altitud?
\begin{opciones}
  \item $\text{Temperatura} = \left(\dfrac{\text{Altitud}}{100}\right) + 10$
  \item $\text{Temperatura} = -\,\text{Altitud} \times 100 + 30$
  \item $\text{Temperatura} = -\left(\dfrac{\text{Altitud}}{100}\right) + 30$
  \item $\text{Temperatura} = \text{Altitud} \times 100 + 10$
\end{opciones}
{\small\color{mmcGris} Icfes, Cuadernillo de preguntas Matemáticas Saber 11.º
(2026), pregunta 8.}
% Ejercicio: icfes-cuadernillo-2026-020 (textual; © Icfes 2026, uso académico)
\pregunta La tabla presenta la información sobre el gasto en publicidad y las
ganancias de una empresa durante los años 2020 a 2022 (datos en millones de
pesos): 2020, gasto en publicidad 200, ganancia obtenida 8.000; 2021, gasto
280, ganancia 10.400; 2022, gasto 250, ganancia 9.500. ¿Cuál es la función que
representa la ganancia obtenida $G$, en millones de pesos, en función del
gasto en publicidad $p$?
\begin{opciones*}(2)
  \item $G(p) = 30p + 2.000$
  \item $G(p) = 10p$
  \item $G(p) = 40p$
  \item $G(p) = 40p - 800$
\end{opciones*}
{\small\color{mmcGris} Icfes, Cuadernillo de preguntas Matemáticas Saber 11.º
(2026), pregunta 20.}
% Ejercicio: icfes-cuadernillo-2026-032 (textual; © Icfes 2026, uso académico)
\pregunta Un trapecio isósceles es un trapecio cuyos lados no paralelos son
congruentes. En un plano cartesiano se dibuja un trapecio isósceles de modo que
el eje $y$ divide al trapecio en dos figuras iguales. Si las coordenadas de dos
de los vértices del trapecio son $(-4, 2)$ y $(-2, 8)$, ¿cuáles son las
coordenadas de los otros dos vértices?
\begin{opciones}
  \item $(8, 2)$ y $(2, 4)$.
  \item $(2, 8)$ y $(4, 2)$.
  \item $(-2, -4)$ y $(-8, -2)$.
  \item $(-4, -2)$ y $(-2, -8)$.
\end{opciones}
{\small\color{mmcGris} Icfes, Cuadernillo de preguntas Matemáticas Saber 11.º
(2026), pregunta 32.}
\end{preguntas}

% ---------------------------------------------------------------------
% 6. Cierre del hilo conductor
% ---------------------------------------------------------------------
\seccion{Cierre: dos números para ubicarse}

De Hiparco a Katherine Johnson, la idea fue siempre la misma: dos números
bastan para ubicar un lugar. Con ellos medimos distancias, encontramos el punto
de encuentro, describimos zonas enteras —todos los puntos a cierta distancia de
una antena, todos los que están igual de lejos de dos estaciones— y trazamos caminos
rectos con una ecuación. Eso es lo que Descartes y Fermat llamaron reducir la
geometría al álgebra. En el siguiente trimestre la circunferencia volverá, junto
con la parábola: son \emph{cónicas}, curvas que estudió Apolonio y que, hacia el
año 400, Hipatia comentaba a sus estudiantes en Alejandría~\cite{dzielska}.

% ---------------------------------------------------------------------
% 7. Evaluación (secciones institucionales)
% ---------------------------------------------------------------------
\matrizevaluacion
  {Reconoce los elementos del plano cartesiano y las fórmulas de distancia,
   punto medio, pendiente y ecuaciones de la recta y de la circunferencia.}
  {Ubica figuras en el plano, calcula distancias y puntos medios, encuentra
   mediatrices, circunferencias y rectas a partir de sus condiciones y las
   grafica.}
  {Justifica sus resultados y los comprueba en el plano antes de darlos por
   buenos.}
  {Comparte sus estrategias y revisa con respeto el trabajo de sus compañeros,
   con argumentos.}

\autoevaluacion

\seguimientodocente

% ---------------------------------------------------------------------
% 8. Referencias
% ---------------------------------------------------------------------
\begin{referencias}
% verificado: https://www.cali.gov.co/informatica/publicaciones/106104/geografia-de-cali/ — «3°27′00″N 76°32′00″O»
\bibitem{alcaldia} Alcaldía de Santiago de Cali. \emph{Geografía de Cali}.
  \url{https://www.cali.gov.co/informatica/publicaciones/106104/geografia-de-cali/}
% verificado: https://fr.wikisource.org/wiki/La_G%C3%A9om%C3%A9trie_(%C3%A9d._1637) — Discours de la méthode plus La Dioptrique, Les Météores et La Géométrie, Leiden, Jan Maire, 1637, pp. 296-413
\bibitem{descartes} Descartes, R. (1637). La Géométrie. En \emph{Discours de
  la méthode} (pp. 296--413). Leiden: Jan Maire.
% verificado: https://bmcr.brynmawr.edu/1995/1995.07.07/ — reseña BMCR de Dzielska, Hypatia of Alexandria, Harvard UP 1995; comentario a las Cónicas de Apolonio en https://mathshistory.st-andrews.ac.uk/Biographies/Hypatia/
\bibitem{dzielska} Dzielska, M. (1995). \emph{Hypatia of Alexandria} (trad.
  F. Lyra). Cambridge, MA: Harvard University Press.
% verificado: https://www.ebsco.com/research-starters/history/hipparchus — «He began the systematic use of longitude and latitude [...] as a means of establishing locations on Earth's surface»
\bibitem{ebsco} EBSCO. \emph{Hipparchus}. Research Starters.
  \url{https://www.ebsco.com/research-starters/history/hipparchus}
% verificado: https://archive.org/details/variaoperamathe00ferm — Varia opera mathematica, Tolosae, apud Johannem Pech, 1679 (ed. Samuel de Fermat)
\bibitem{fermat} Fermat, P. de (1679). \emph{Varia opera mathematica}.
  Toulouse: J. Pech.
% verificado: recursos/matematicas/icfes/09-Marzo_Cuadernillo-de-Preguntas-Matematicas-Saber-11-2026.pdf (archivo local) — pregunta 32, clave B
\bibitem{icfes} Icfes (2026). \emph{Cuadernillo de preguntas: Matemáticas,
  Saber 11.º}. Bogotá: Instituto Colombiano para la Evaluación de la Educación.
% verificado: https://mathshistory.st-andrews.ac.uk/Biographies/Descartes/ — La Géométrie, uno de los tres apéndices del tratado publicado en Leiden en 1637, «by far the most important part»
\bibitem{mactutordescartes} O'Connor, J. J. y Robertson, E. F. \emph{René
  Descartes}. MacTutor History of Mathematics, Universidad de St Andrews.
% verificado: https://mathshistory.st-andrews.ac.uk/Biographies/Hipparchus/ — nacido en 190 a. C. en Nicea, muerto hacia 120 a. C.
\bibitem{mactutorhiparco} O'Connor, J. J. y Robertson, E. F. \emph{Hipparchus
  of Rhodes}. MacTutor History of Mathematics, Universidad de St Andrews.
% verificado: https://www.encyclopedia.com/people/science-and-technology/mathematics-biographies/pierre-de-fermat — M. S. Mahoney, Complete Dictionary of Scientific Biography: Isagoge compuesta antes de la primavera de 1636, independiente de Descartes; Dx = Ey (recta)
\bibitem{mahoney} Mahoney, M. S. Fermat, Pierre de. En \emph{Complete
  Dictionary of Scientific Biography}. Charles Scribner's Sons (Encyclopedia.com).
% verificado: https://gblumen.mineducacion.gov.co/cgi-bin/koha/opac-detail.pl?biblionumber=4785 — MEN 2016, ISBN 978-958-691-913-5
\bibitem{mendba} Ministerio de Educación Nacional (2016). \emph{Derechos
  Básicos de Aprendizaje V.2: Matemáticas}. Bogotá: MEN.
% verificado: https://www.nasa.gov/centers-and-facilities/langley/katherine-johnson-biography/ — 1918-2020; Shepard (1961), Glenn (1962), Medalla Presidencial de la Libertad (2015)
\bibitem{nasajohnson} NASA. \emph{Katherine Johnson Biography}.
  \url{https://www.nasa.gov/centers-and-facilities/langley/katherine-johnson-biography/}
% verificado: https://oceanservice.noaa.gov/facts/latitude.html — «Each degree of latitude covers about 111 kilometers on the Earth's surface»
\bibitem{noaa} NOAA National Ocean Service. \emph{What is latitude?}
  \url{https://oceanservice.noaa.gov/facts/latitude.html}
% verificado: https://catalog.hathitrust.org/Record/011445470 — NASA TN D-233, Langley, 1960; Skopinski y Johnson
\bibitem{skopinski} Skopinski, T. H. y Johnson, K. G. (1960).
  \emph{Determination of Azimuth Angle at Burnout for Placing a Satellite Over
  a Selected Earth Position} (NASA TN D-233). Washington: NASA.
\end{referencias}

\end{document}
